\documentclass{article}
\usepackage{amssymb}
\usepackage{amsmath}
 
\title{Maximum Number of Electron Diffraction Rings}
\author{Joe Howlett}
\date{August 2026}

\begin{document}
\maketitle
\begin{abstract}
\noindent In a model of the electron diffraction experiment by Lester Germer and Clinton Davisson the number of electron diffraction rings can be calculated using Bragg's law and the de Broglie wavelength equation.
\end{abstract}

\section*{Maximum value of n}
\noindent Where $n$ is the number of rings, $h$ is planks constant, $m$ is mass of electron, $e$ is charge of electron, $V$ is potential difference of $5kV$ and $d$ is the atomic spacing in layers of graphite atoms. The C--C bond length in graphite is approximately
$0.142\,\mathrm{nm}$ (Sun et al., 2009). Therefore $d_{1}$ is $\frac{\sqrt{3}}{2}\times {0.142\times 10^{-9}} = 0.123\times10^{-9}${\scriptsize\text{(3sf)}} and $d_{2}$ is $\frac{3}{2}\times {0.142\times 10^{-9}} = 0.213\times10^{-9}$.

\vspace{1em}

\noindent n is in the relationship:
\begin{align*}
\sin\frac{\phi}{2}=\frac{nh}{2d\sqrt{2 m e V}} \\
\intertext{Isolating n:}
\frac{\sin\frac{\phi}{2}{2}d\sqrt{2 m e V}}{h}=n \\
\intertext{Due to the maximum value of sin being $1$:}
\frac{{2}d\sqrt{{2}meV}}{{h}}\ge{n}& \\
\therefore \frac{{2}d\sqrt{{2}meV}}{{h}} ={n}&_{\text{max}}
\end{align*}

\noindent When $d$ is equal to $0.123\times10^{-9}: {n}_{\text{max}} = 14.2$ {\scriptsize\text{(3sf)}}

\noindent When $d$ is equal to $0.213\times10^{-9}: {n}_{\text{max}} = 24.5$ {\scriptsize\text{(3sf)}}

\vspace{1em}

\noindent As n is the maximum number of rings, the values should be truncated to integers.
\begin{align*}
d = 0.123\times10^{-9},\ n= 14 \\
d = 0.213\times10^{-9},\ n= 24
\end{align*}

\end{document}
